\documentclass[11pt,a4paper]{article}
\usepackage[margin=21mm]{geometry}
\usepackage{amsmath,amssymb,booktabs,array}
\usepackage[hidelinks]{hyperref}
\setlength{\parindent}{0pt}\setlength{\parskip}{5pt}
\raggedright\emergencystretch=3em
\begin{document}
\begin{center}{\Large D2: node consistency, spatial state and the next clock}\\[4pt]
9 October 2026 --- methodology-only qualification\end{center}
\textbf{Status.} D2 is approved. P.A/T.A/N.A/G1/E1/K.A remain fixed.
This note qualifies a static node block and explicit field/reset diagnostics.
The complete endogenous field, clearance and control/gradient claims are not
proved. The regular-clock field contract Q.J.F.R remains open.

\section*{1. Definitions and the selected static map}
Population is continuous vehicle-equivalents (eq); distance is m and time s.
For the diagnostic triangular first-order flux,
\[
 F(\rho)=\min\{u\rho,w(\kappa-\rho)\},\quad
 u=30,\ w=5,\ \kappa=3/20,\quad
 \rho_c=3/140,\quad Q=9/14\ \mathrm{eq/s}.
\]
At the public request clock, $S_a$ is a sampled sending \emph{amount} and
$\theta_{ab}$ is its direction fraction. The original amount allowance and
competition weight are $d_{ab}=\mu_{ab}=\theta_{ab}S_a$ (eq).
An active port has unspent allowance. During physical evolution,
\[
 b_{ab}(t)=d_{ab}-\int_{t_j}^{t}f_{ab}(r)\,dr\ge0.
\]
Here $f_{ab}$ is a rate (eq/s); fulfilled ports remain inactive until the
next public refresh. $D_a$ is current physical donor sending (eq/s),
$R_b$ current receiving, and $Q_a$ donor capacity. D2 imposes
\[
 f_{ab}\ge0,\qquad \sum_b f_{ab}\le D_a\le Q_a,\qquad
 \sum_a f_{ab}\le R_b,
\]
without the D1 inequality $f_{ab}\le\theta_{ab}Q_a$.

For the ramp target $A$, let $D_r$ be available ramp sending and
$p\in[0,1]$ its fixed priority. The resource available to the mainline
subgroup and the completed ramp award are
\[
 C_A=R_A-\min\{pR_A,D_r\},\qquad
 f_r=\min\{D_r,R_A-m_A\},\quad m_A=\sum_a f_{aA}.
\]
Thus $m_A+f_r\le R_A$, with unused mainline space returned to the ramp.
The mainline allocation uses $C_A$ and the other receivers' supplies.
Starting from zero, increase rates of all live ports proportionally to their
fixed positive $\mu_{ab}$. Stop at the first saturated row or column,
freeze every incident live port at all ties, and repeat. Each pass freezes
at least one port, so the static calculation terminates in at most the
number of active ports. Its water level has units $1/\mathrm{s}$ and is
\emph{not physical execution time}. No FIFO/throughput theorem is imported.

\newpage
\section*{2. Why own boundary-state changes do not alter this allocation}
\textbf{Incoming demand lift.} Let $q_a=\sum_b f_{ab}<D_a$ after allocation.
The row constraint was slack throughout all monotone increments. Replacing
$D_a$ by $Q_a$ therefore cannot affect any first saturation or freeze.
Induction over the finite passes leaves every $f_{ab}$ unchanged.
For the ramp, if $f_r<D_r$, then $f_r=R_A-m_A$. This cannot occur when
$D_r<pR_A$, since then $m_A\le R_A-D_r$. Hence $D_r\ge pR_A$,
$C_A=(1-p)R_A$, and raising $D_r$ to capacity leaves both $C_A$ and
$f_r$ unchanged. Multiple strict demand lifts may be made together.

\textbf{Unfilled receiving lift.} A slack ordinary column cannot have caused
a freeze; raising its bound leaves the same increment sequence. If the ramp
target is slack, $m_A+f_r<R_A$, then $f_r=D_r$ and
\[
 m_A<R_A-D_r\le C_A.
\]
The mainline resource is slack too. Since
$C_A(R_A)=R_A-\min\{pR_A,D_r\}$ is nondecreasing, raising $R_A$ cannot
change the mainline allocation or the fully served ramp award.

In triangular KW, a strictly restricted sender admits a congested boundary
trace with demand $Q_a$; an unfilled receiver admits a free inlet trace with
supply $Q_b$. These lifts are therefore the relevant stationary-boundary
consistency checks. Original weights and active requests must remain fixed.
An amount exhaustion, new clock or restriction encounter is a genuine
physical event, not covered by this fixed-mode proof.

\textbf{Supplementary exact checks.} The declared two-donor/two-recipient/ramp
layout was checked on 5,832 rational configurations: zero resources,
asymmetric fixed weights, two ramp priorities and exhausted fork ports.
All 8,292 strict demand lifts and 3,716 slack supply lifts leave the allocation
unchanged. Enumeration supplements the argument; it is not whole-model validation.

\section*{3. What the regular clock still needs}
Restriction anchors constrain the mass between selected points, not every
ordinary density piece. If local count has gauge
\[
 N(t,x)=n(0)+\int_0^t q_{\rm in}(r)dr-\int_0^x\rho(t,y)dy,
\]
two fields can share $N(t,0)$, $N(t,L)$ and all record anchors while having
different density near the exit. Replacing both by a linear count interpolation
preserves selected counts but is not an identity map on traffic fields.

\newpage
\section*{4. Exact early/late profiles under the same D2 request data}
Let $L=30$ m, $T=1$ s, the through receiver be blocked and the lateral
receiver empty/free. There are no donor records in this diagnostic.
Both initial fields contain $n_0=3\kappa=9/20$ eq and use the same already
sampled $\theta_L=1/4$, giving $d_L=9/80$ eq. They have identical fixed-end
counts and any shared records outside this donor can be held identical.
They are admissible field/request inputs, not proven equal-history policies.

\textbf{Early: jam on $[27,30]$.} The exit can send $Q$ immediately.
The allowance is exhausted at
\[
 t_e=d_L/Q=7/40,\qquad M_e=Q t_e=9/80=0.1125\ \mathrm{eq}.
\]
The initial release contact travels left at $-w$, reaching 27 m at $3/5$ s.
Afterward the void/critical rear moves at $u$. The quota-stop contact from
the exit also moves left at $-w$. They meet at
\[
 27+30(t-3/5)=30-5(t-7/40),\qquad t=5/8,\quad x=111/4.
\]
The remaining field at $T$ is stationary jam on $[111/4,30]$ and void behind.
The lateral recipient contains critical density on $[99/4,30]$:
\[
 n_e(T)=\kappa(30-111/4)=27/80=n_0-M_e,
 \qquad \rho_c(30-99/4)=9/80=M_e.
\]

\textbf{Late: jam on $[0,3]$.} The leading void/critical front travels at
$u$ and reaches the exit at $t_a=(30-3)/30=9/10$. Hence
\[
 M_l=Q(1-t_a)=9/140\simeq0.064286<d_L.
\]
The release contact reaches the inlet at $3/5$, after which the free rear
travels to $x=30(1-3/5)=12$ at $T$. The completed donor is critical on
$[12,30]$; its recipient is critical on $[0,3]$:
\[
 n_l(T)=18\rho_c=27/70=n_0-M_l,\qquad 3\rho_c=9/140=M_l.
\]
Both conserve the common population exactly. Their accepted mass differs by
$M_e-M_l=27/560$ eq. Anchor-only uniform reconstruction instead maps both
initial donors to $\rho=n_0/L=3/200$. It cannot be equivalent to both fields.
No positive traffic has exited the diagnostic recipient before $T$; the
one-step total population integral is identical. Thus this is a field-state
non-equivalence result, not a full-TST benefit or mesoscopic-necessity witness.

\newpage
\section*{5. Completed-state E1 is compatible with the next input}
E1 reads completed transfers and creates zero-population records at the
declared ports. Every record at the same lane/time/point attaches to the same
$N(t,x)$; different signatures remain distinct. It does not alter $\rho$,
population, boundary counts or exits. For two source-free records,
\[
 x_1\le x_2,\quad h_i=N(t,x_i),\quad
 h_1-h_2=\int_{x_1}^{x_2}\rho(t,x)dx,
 \quad x_2-x_1\ge\frac{h_1-h_2}{\kappa}.
\]
Unequal counts cannot coincide at finite jam density. Equal counts enclose
void and admit the prior local coincident fan without new ranks or compression.
This is conditional reset compatibility, not proof of all later contacts.

For the earlier common one-step A-section certificate,
\[
 \rho=9/100\text{ on }[0,13/3],\quad n=39/100,
 \quad N(1,0)=39/100,\quad N(1,13/3)=0.
\]
Both new inlet records receive $39/100$, while the old record keeps zero.
This gauge includes initial mass $9/100$ and actual entry $3/10$; it is not
a claim that $39/100$ new mass entered. Their initial separation is $13/3$ m,
above $(39/100)/(3/20)=13/5=2.6$ m. Own queue speed is $10/3$ m/s;
birth speeds are capped by that same own field, with distinct signatures retained.

\textbf{Remaining reference decision.} R.A carries the completed full density
or equivalent count profile to the next interval. R.B declares a numerical
mass-and-anchor reconstruction and separately qualifies its field/timing error.
Both must retain all records and selected anchors. The early/late pair refutes
calling the R.B projection an exact identity; it does not refute a qualified
finite-grid approximation or macro modeling. No representation is adopted here.

\textbf{Whole-model limit.} Static node consistency and E1 compatibility do
not prove endogenous motion/port existence for every admitted encounter,
indefinite continuation, clearance, intervention-dependent full TST or useful
gradients. A coherent finite numerical construction at a declared accuracy
could qualify the chain without a universal continuum theorem. Its contract
must be bound before using it as a competing baseline.
\end{document}
